% GENERATED FILE. Edit canonical lessons, then run npm run generate:books.
% canonical-source-hash: fnv1a64:6d6c2007eaa1fb56
% canonical-lessons: FA-C455-shahrbazi, FA-C455-park-e-jangali, FA-C455-namayeshgah, FA-C455-farhangsara, FA-C455-tatil-shodan

\chapter{Out and about}
\label{ch:fa-out-and-about}
\hlchaptermodality{455}

\begin{chapteropening}
\textbf{By the end of this chapter:} \emph{I can talk about places to go in town: the funfair, a forest park, an exhibition, a cultural centre and closing time.}

Complete the last lesson of chapter 455: I can talk about places to go in town: the funfair, a forest park, an exhibition, a cultural centre and closing time.
\end{chapteropening}

\section[\texorpdfstring{\fa{شهربازی} (shahrbâzi)}{شهربازی (shahrbâzi)}]{\fa{شهربازی} (shahrbâzi) — an amusement park, a funfair}

\label{lesson:FA-C455-shahrbazi}



\begin{quote}
Before the new one: say the Persian for a hairdresser's, then the Persian for a watchmaker.
\end{quote}

\subsection*{You'll want to know: \fa{شهربازی}}
\textbf{\fa{شهربازی}} (\emph{shahrbâzi}) — \textquotedblleft{}an amusement park, a funfair\textquotedblright{}.

\begin{quote}
\fa{بچه} \fa{ها} \fa{روز} \fa{جمعه} \fa{به} \fa{شهربازی} \fa{رفتند}.
\end{quote}

\begin{quote}
\emph{Bachche-hâ ruz-e jom'e be shahrbâzi raftand.}
\end{quote}

\begin{quote}
The children went to the funfair on Friday.
\end{quote}

Literally 'play-town'.

\begin{cousinweb}
\emph{shahrbâzi} is a native compound of \emph{shahr} 'town', from Middle Persian \emph{šahr} (EtymDB), and \emph{bâzi} 'play', usually traced to Middle Persian \emph{wāzīg}.
\end{cousinweb}

\begin{grammarlens}[title={What you've built}]
One more word for telling, explaining and describing at B1.
\end{grammarlens}

\subsection*{Your turn}
\begin{itemize}
\raggedright
  \item \emph{Say it:} \emph{shahrbâzi}
  \item \emph{Say it:} \emph{shahrbâzi}, once more
  \item \emph{Say it:} a sentence of your own with \emph{shahrbâzi}
  \item \emph{Recall:} say the Persian for a hairdresser's, then the Persian for a watchmaker, then say \emph{shahrbâzi} again
\end{itemize}

\begin{tcolorbox}[breakable,colback=teal!4,colframe=teal!35!black,title={Before you move on}]
What does \fa{شهربازی} mean? (\textquotedblleft{}An amusement park, a funfair\textquotedblright{}.) Say it once more.
\end{tcolorbox}
\section[\texorpdfstring{\fa{پارک} \fa{جنگلی} (pârk-e jangali)}{پارک جنگلی (pârk-e jangali)}]{\fa{پارک} \fa{جنگلی} (pârk-e jangali) — a forest park, a wooded park}

\label{lesson:FA-C455-park-e-jangali}



\begin{quote}
Before the new one: say the Persian for a watchmaker, then the Persian for an amusement park.
\end{quote}

\subsection*{You'll want to know: \fa{پارک} \fa{جنگلی}}
\textbf{\fa{پارک} \fa{جنگلی}} (\emph{pârk-e jangali}) — \textquotedblleft{}a forest park, a wooded park\textquotedblright{}.

\begin{quote}
\fa{جمعه} \fa{ها} \fa{مردم} \fa{در} \fa{پارک} \fa{جنگلی} \fa{کباب} \fa{درست} \fa{میکنند}.
\end{quote}

\begin{quote}
\emph{Jom'e-hâ mardom dar pârk-e jangali kabâb dorost mikonand.}
\end{quote}

\begin{quote}
On Fridays people make kebabs in the forest park.
\end{quote}

Many Iranian towns have a \emph{pârk-e jangali} on their edge for picnics.

\begin{cousinweb}
\emph{pârk-e jangali} is built from \emph{pârk}, a loan usually traced to French \emph{parc}, and \emph{jangali} 'wooded', from \emph{jangal} 'forest', usually traced to Sanskrit \emph{jaṅgala}.
\end{cousinweb}

\begin{grammarlens}[title={What you've built}]
One more word for telling, explaining and describing at B1.
\end{grammarlens}

\subsection*{Your turn}
\begin{itemize}
\raggedright
  \item \emph{Say it:} \emph{pârk-e jangali}
  \item \emph{Say it:} \emph{pârk-e jangali}, once more
  \item \emph{Say it:} a sentence of your own with \emph{pârk-e jangali}
  \item \emph{Recall:} say the Persian for a watchmaker, then the Persian for an amusement park, then say \emph{pârk-e jangali} again
\end{itemize}

\begin{tcolorbox}[breakable,colback=teal!4,colframe=teal!35!black,title={Before you move on}]
What does \fa{پارک} \fa{جنگلی} mean? (\textquotedblleft{}A forest park, a wooded park\textquotedblright{}.) Say it once more.
\end{tcolorbox}
\section[\texorpdfstring{\fa{نمایشگاه} (namâyeshgâh)}{نمایشگاه (namâyeshgâh)}]{\fa{نمایشگاه} (namâyeshgâh) — an exhibition; a trade fair}

\label{lesson:FA-C455-namayeshgah}



\begin{quote}
Before the new one: say the Persian for an amusement park, then the Persian for a forest park.
\end{quote}

\subsection*{You'll want to know: \fa{نمایشگاه}}
\textbf{\fa{نمایشگاه}} (\emph{namâyeshgâh}) — \textquotedblleft{}an exhibition; a trade fair\textquotedblright{}.

\begin{quote}
\fa{هفته} \fa{پیش} \fa{به} \fa{نمایشگاه} \fa{کتاب} \fa{رفتیم}.
\end{quote}

\begin{quote}
\emph{Hafte-ye pish be namâyeshgâh-e ketâb raftim.}
\end{quote}

\begin{quote}
Last week we went to the book fair.
\end{quote}

\emph{namâyesh} 'show' + \emph{gâh} 'place'.

\begin{cousinweb}
\emph{namâyeshgâh} is a native compound of \emph{namâyesh} 'a showing', built on \emph{nemudan} 'to show', and \emph{gâh} 'place'.
\end{cousinweb}

\begin{grammarlens}[title={What you've built}]
One more word for telling, explaining and describing at B1.
\end{grammarlens}

\subsection*{Your turn}
\begin{itemize}
\raggedright
  \item \emph{Say it:} \emph{namâyeshgâh}
  \item \emph{Say it:} \emph{namâyeshgâh}, once more
  \item \emph{Say it:} a sentence of your own with \emph{namâyeshgâh}
  \item \emph{Recall:} say the Persian for an amusement park, then the Persian for a forest park, then say \emph{namâyeshgâh} again
\end{itemize}

\begin{tcolorbox}[breakable,colback=teal!4,colframe=teal!35!black,title={Before you move on}]
What does \fa{نمایشگاه} mean? (\textquotedblleft{}An exhibition; a trade fair\textquotedblright{}.) Say it once more.
\end{tcolorbox}
\section[\texorpdfstring{\fa{فرهنگسرا} (farhangsarâ)}{فرهنگسرا (farhangsarâ)}]{\fa{فرهنگسرا} (farhangsarâ) — a cultural centre, an arts centre}

\label{lesson:FA-C455-farhangsara}



\begin{quote}
Before the new one: say the Persian for a forest park, then the Persian for an exhibition.
\end{quote}

\subsection*{You'll want to know: \fa{فرهنگسرا}}
\textbf{\fa{فرهنگسرا}} (\emph{farhangsarâ}) — \textquotedblleft{}a cultural centre, an arts centre\textquotedblright{}.

\begin{quote}
\fa{در} \fa{فرهنگسرای} \fa{محله} \fa{کلاس} \fa{موسیقی} \fa{هست}.
\end{quote}

\begin{quote}
\emph{Dar farhangsarâ-ye mahalle kelâs-e musiqi hast.}
\end{quote}

\begin{quote}
There are music classes at the local cultural centre.
\end{quote}

\emph{farhang} 'culture' + \emph{sarâ} 'house'.

\begin{cousinweb}
\emph{farhangsarâ} is a native compound of \emph{farhang} 'culture', from Middle Persian \emph{frahang} (EtymDB), and \emph{sarâ} 'house, hall'.
\end{cousinweb}

\begin{grammarlens}[title={What you've built}]
One more word for telling, explaining and describing at B1.
\end{grammarlens}

\subsection*{Your turn}
\begin{itemize}
\raggedright
  \item \emph{Say it:} \emph{farhangsarâ}
  \item \emph{Say it:} \emph{farhangsarâ}, once more
  \item \emph{Say it:} a sentence of your own with \emph{farhangsarâ}
  \item \emph{Recall:} say the Persian for a forest park, then the Persian for an exhibition, then say \emph{farhangsarâ} again
\end{itemize}

\begin{tcolorbox}[breakable,colback=teal!4,colframe=teal!35!black,title={Before you move on}]
What does \fa{فرهنگسرا} mean? (\textquotedblleft{}A cultural centre, an arts centre\textquotedblright{}.) Say it once more.
\end{tcolorbox}
\section[\texorpdfstring{\fa{تعطیل} \fa{شدن} (ta'til shodan)}{تعطیل شدن (ta'til shodan)}]{\fa{تعطیل} \fa{شدن} (ta'til shodan) — to close (for the day or a holiday), to shut}

\label{lesson:FA-C455-tatil-shodan}



\begin{quote}
Before the new one: say the Persian for an exhibition, then the Persian for a cultural centre.
\end{quote}

\subsection*{You'll want to know: \fa{تعطیل} \fa{شدن}}
\textbf{\fa{تعطیل} \fa{شدن}} (\emph{ta'til shodan}) — \textquotedblleft{}to close (for the day or a holiday), to shut\textquotedblright{}.

\begin{quote}
\fa{بازار} \fa{ساعت} \fa{هشت} \fa{شب} \fa{تعطیل} \fa{شد}.
\end{quote}

\begin{quote}
\emph{Bâzâr sâ'at-e hasht-e shab ta'til shod.}
\end{quote}

\begin{quote}
The bazaar closed at eight in the evening.
\end{quote}

\emph{ta'til} 'closed, a holiday' + \emph{shodan}. \emph{Ruz-e ta'til} — a day off.

\begin{cousinweb}
\emph{ta'til shodan} is built from \emph{ta'til} 'closure, a holiday', borrowed from Arabic \emph{taʿṭīl} 'suspension, a break' (EtymDB), and the native verb \emph{shodan} 'to become'.
\end{cousinweb}

\begin{grammarlens}[title={What you've built}]
One more word for telling, explaining and describing at B1.
\end{grammarlens}

\subsection*{Your turn}
\begin{itemize}
\raggedright
  \item \emph{Say it:} \emph{ta'til shodan}
  \item \emph{Say it:} \emph{ta'til shodan}, once more
  \item \emph{Say it:} a sentence of your own with \emph{ta'til shodan}
  \item \emph{Recall:} say the Persian for an exhibition, then the Persian for a cultural centre, then say \emph{ta'til shodan} again
\end{itemize}

\begin{tcolorbox}[breakable,colback=teal!4,colframe=teal!35!black,title={Before you move on}]
What does \fa{تعطیل} \fa{شدن} mean? (\textquotedblleft{}To close (for the day or a holiday), to shut\textquotedblright{}.) Say it once more.
\end{tcolorbox}
